% ======================================================================
% 第 III 节：规范场的格点量子化
% ======================================================================
\section{规范场的格点量子化}\label{sec:lattice}

\subsection{格点上的经典作用量}\label{sec:classical}

本节将把规范场作用量（拉格朗日量的时空积分）定义在欧几里得时空间距为 $a$ 的
离散格点上。最简单的做法是考虑连续作用量，用有限差分近似替代导数，并把时空积分
换成对格点的求和。然而这样得到的作用量对非零 $a$ 不是规范不变的。考虑到重正化
的不确定性，这很可能意味着量子化理论在 $a \to 0$ 极限下（如果该极限存在）仍然
缺乏规范不变性。另一种做法是为格点理论表述规范不变性，并调整作用量使之对任意
$a$ 规范不变。本文采用后一种做法。

为方便起见，将详细讨论一个带电荷的狄拉克场 $\psi$ 与单一规范场 $A_{\mu}$ 耦合
的情形。对非阿贝尔规范群的推广将在后文指出。

在格点上，场量为 $\psi_{n}$、$\bar\psi_{n}$ 和 $\thl{\mu\nu}(n)$，其中 $n$ 是
分量均为整数的四矢量，指向简单超立方格点上的点。

（暂不考虑规范场的）狄拉克场的简单格点作用量具有如下形式
\begin{equation}
A_{D} = -a^{3}\sum_{n} m_{b}\bar\psi_{n}\psi_{n}
        + a^{3}\sum_{n,\mu} \bar\psi_{n}\gamma_{\mu}
          \frac{\psi_{n+\hat\mu}-\psi_{n-\hat\mu}}{2a},
\label{eq:3.1}
\end{equation}
其中 $m_{b}$ 是裸质量；$\hat\mu$ 是沿 $\mu$ 轴的单位矢量；$a^{3}\sum_{n}$ 替代
连续理论的时空积分，$(\psi_{n+\hat\mu}-\psi_{n-\hat\mu})/2a$ 替代
$\partial_{\mu}\psi$。由于采用欧几里得度规，没有总的 $i$ 因子。格点上的规范变换
可以定义如下。
\begin{equation}
\psi_{n} \to e^{i\Lambda_{n}}\psi_{n},\qquad
\bar\psi_{n} \to \bar\psi_{n}e^{-i\Lambda_{n}},
\label{eq:3.2}
\end{equation}
\begin{equation}
\thl{\mu\nu}(n) \to \thl{\mu\nu}(n) - \frac{\gamma_{n+\hat\mu}-\gamma_{n}}{a},
\label{eq:3.3}
\end{equation}
\stepcounter{equation}%   % 论文中没有 (3.4)
其中 $g$ 是耦合常数，$\gamma_{n}$ 是任意的。项 $\bar\psi_{n}\psi_{n+\hat\mu}$
和 $\bar\psi_{n+\hat\mu}\psi_{n}$ 对此变换不变；相应的规范不变表达式是
$\bar\psi_{n}\psi_{n+\hat\mu}\exp(iga\,\thl{\mu\nu}(n))$ 和
$\bar\psi_{n}\psi_{n+\hat\mu}\exp(-iga\,\thl{\mu\nu}(n))$。因此，$A_{D}$ 的
规范不变形式是
\begin{equation}
\begin{split}
A_{D} = {} & -a^{3}\sum_{n} m_{b}\bar\psi_{n}\psi_{n}\\
           & + \frac{a^{2}}{2}\sum_{n,\mu}
             \Bigl[\bar\psi_{n}\gamma_{\mu}\psi_{n+\hat\mu}
             e^{iga\,\thl{\mu\nu}(n)}
             -\bar\psi_{n+\hat\mu}\gamma_{\mu}\psi_{n}
             e^{-iga\,\thl{\mu\nu}(n)}\Bigr].
\end{split}
\label{eq:3.5}
\end{equation}

在作用量 $A_{D}$ 中，场 $\thl{\mu\nu}$ 表现得像一个角变量：$A_{D}$ 对
$\thl{\mu\nu}$ 以 $2\pi$ 为周期。

自由规范场作用量将被定义为保持这一性质。这并不意味着在连续极限下
$\thl{\mu\nu}$ 是角变量。由于关系式~\eqref{eq:3.6}，当 $a \to 0$ 时
$\thl{\mu\nu}$ 变成无穷小角变量 $A_{\mu}(x)$；这样的变量取值范围为
$-\infty < A_{\mu}(x) < \infty$，没有任何周期性。
\begin{equation}
\thl{\mu\nu}(n) = a\,g\,A_{\mu}(na) + \order{a^{2}}.
\label{eq:3.6}
\end{equation}

$F_{\mu\nu}$ 的一个规范不变的格点近似是
\begin{equation}
F_{\mu\nu}(n) =
\frac{\thl{\mu\nu}(n)-\thl{\nu\mu}(n)-\thl{\mu\nu}(n+\hat\nu)
+\thl{\nu\mu}(n+\hat\mu)}{a}.
\label{eq:3.7}
\end{equation}

定义 $\thl{\mu\nu}$ 的一个重定标形式是方便的：
\begin{equation}
f_{\mu\nu}(n) = \thl{\mu\nu}(n)+\thl{\nu\mu}(n+\hat\mu)
-\thl{\mu\nu}(n+\hat\nu)-\thl{\nu\mu}(n).
\label{eq:3.8}
\end{equation}

保持周期性的规范场的简单格点作用量是
\begin{equation}
A_{G} = \frac{1}{2g^{2}}\sum_{n,\mu\nu}
        \Bigl[1-\tfrac{1}{2}\bigl(e^{if_{\mu\nu}(n)}+e^{-if_{\mu\nu}(n)}\bigr)\Bigr].
\label{eq:3.9}
\end{equation}

在连续极限下，由于 $f_{\mu\nu}$ 定义中的因子 $a$，$f_{\mu\nu} \to 0$。因此对
小的 $a$ 可以写
\[
A_{G} = \frac{1}{2g^{2}}\sum_{n,\mu\nu}
        \Bigl[\tfrac{1}{4}\bigl(f_{\mu\nu}(n)\bigr)^{2}+\cdots\Bigr],
\]
常数项无关紧要。$f_{\mu\nu}$ 中的线性项为零，因为 $f_{\mu\nu}$ 对指标 $\mu$
和 $\nu$ 是反对称的。二次项给出
\stepcounter{equation}%   % 论文中没有 (3.10)
\begin{equation}
A_{G} = -\frac{1}{2g^{2}}\,a^{4}\sum_{n,\mu\nu}\bigl(F_{\mu\nu}(n)\bigr)^{2},
\label{eq:3.11}
\end{equation}
这正是格点近似下通常的规范场作用量。涉及 $f_{\mu\nu}^{3}$、$f_{\mu\nu}^{4}$
等各项即使在提出因子 $a^{4}$ 把 $\sum_{n}$ 化为积分之后，在 $a\to0$ 时也
全部消失。

现在可以写出完整作用量
\begin{equation}
\begin{split}
A = {} & -\epsilon\sum_{n}\bar\psi_{n}\psi_{n}
         + \kappa\sum_{n,\mu}\Bigl[
           \bar\psi_{n+\hat\mu}\gamma_{\mu}\psi_{n}e^{i\thl{\mu\nu}(n)}
           -\bar\psi_{n}\gamma_{\mu}\psi_{n+\hat\mu}e^{-i\thl{\mu\nu}(n)}\Bigr]\\
       & + \frac{1}{2g^{2}}\sum_{n,\mu\nu}
         \Bigl[1-\tfrac{1}{2}\bigl(e^{if_{\mu\nu}(n)}+e^{-if_{\mu\nu}(n)}\bigr)\Bigr],
\end{split}
\label{eq:3.12}
\end{equation}
其中 $\epsilon = m_{0}a^{4}$，$\kappa = a^{2}/2$。这一作用量在 $a \to 0$ 时
回到通常的连续作用量；对有限的 $a$ 它规范不变且对规范场具有周期性。不过要注意，
连续极限是一个经典极限，其中格点变量 $\psi_{n}$、$\bar\psi_{n}$ 和
$\thl{\mu\nu}$ 趋于连续函数 $\psi(x)$、$\bar\psi(x)$ 和 $A_{\mu}(x)$，
其中 $x=na$。量子化理论的连续极限由于重正化问题而更难讨论。

\subsection{量子化}\label{sec:quantization}

这里最关心的问题是规范场的量子化。因此先单独量子化规范场。稍后再讨论 $\psi$ 的
量子化。本节末尾还将描述对非阿贝尔规范理论的推广。

格点规范理论的量子化分两步进行。第一步是定义格点版本的欧几里得真空期望值，从
格点版本的费曼路径积分出发。第二步是定义格点上的量子理论，从而引入实时间变量并
定义粒子态和散射振幅。两种情况都由于格点提供紫外截断而没有洛伦兹不变性。洛伦兹
不变性只能在 $a$（格点间距）$\to 0$ 的极限下（如果真能达到）实现，而实际上这是
一个难以评估的极限。

正如第一节所讨论的，我们希望计算 $\exp[ig\oint A_{\mu}(x)\,ds^{\mu}]$ 以规范场
作用量为权重的规范场平均。在格点上，线积分变成对格点上一条闭合路径 $P$ 的和
（例如见图~\ref{fig:8}）。该和具有形式 $ig\sum_{P}(\pm)\thl{\mu\nu}$，其中
特定的 $\thl{\mu\nu}$ 出现在求和中，当路径连接格点 $n$ 与 $n+\hat\mu$ 时取
$+\thl{\mu\nu}$（当路径从 $n+\hat\mu$ 走向 $n$ 时取 $-\thl{\mu\nu}$）。

\begin{figure}[t]
  \centering
  \includegraphics[width=0.4\textwidth]{fig8}
  \refstepcounter{figure}%
  \caption*{图 8. 格点路径 $P$ 的一个例子。}
  \label{fig:8}
\end{figure}

在格点上，对一切规范场求平均涉及对所有 $n$ 和 $\mu$ 的所有 $\thl{\mu\nu}$ 值
积分。通常积分范围为无限区间：$-\infty < \thl{\mu\nu} < \infty$，但由于
$\thl{\mu\nu}$ 具有周期性，积分超过一个周期毫无意义。因此规范场平均的格点版本是
\begin{equation}
I_{\theta}(P) = \prod_{n,\mu\nu}\int\ddth{\mu\nu}
  \exp\!\Bigl[i\sum_{P}(\pm)\thl{\mu\nu}
    +\frac{1}{2g^{2}}\sum_{n,\mu\nu}
      \bigl(e^{if_{\mu\nu}(n)}+e^{-if_{\mu\nu}(n)}\bigr)\Bigr].
\label{eq:3.13}
\end{equation}
\begin{equation}
Z_{\theta} = \prod_{n,\mu\nu}\int\ddth{\mu\nu}
  \exp\!\Bigl[\frac{1}{2g^{2}}\sum_{n,\mu\nu}
    \bigl(e^{if_{\mu\nu}(n)}+e^{-if_{\mu\nu}(n)}\bigr)\Bigr].
\label{eq:3.14}
\end{equation}

注意没有在作用量中加入规范固定项。$\thl{\mu\nu}$ 的有限取值范围使得规范固定项
没有必要。在连续规范理论中，$A_{\mu}(x)$ 的取值范围为
$[-\infty < A_{\mu}(x) < \infty]$，规范固定项对于使泛函积分收敛必不可少。原因在于
所有可能规范变换在路径积分空间中所张的体积是无穷大的；规范固定项为该体积提供
收敛因子。在格点理论中，如果格点本身有界，积分总体积就是有限的，无需收敛因子。
对于无穷大范围的格点，由于积分次数无穷多（换言之格点体积无穷大）会出现发散，
但这些发散通常通过对 $Z$ 的除法（这一除法等价于在微扰论中去除所有真空圈）消除。

可以类似地定义更常规的真空期望值。例如，可以定义传播子。在没有规范固定项时，
直接为规范场 $\thl{\mu\nu}$ 定义传播子很不方便；作为替代，可以定义规范不变的
传播子为
\begin{equation}
D_{\mu\nu}(n) = Z_{\theta}^{-1}\prod_{m,\rho}\int\ddth{\rho\sigma}
  e^{i\thl{\mu\nu}(n)}e^{-i\thl{\mu\nu}(0)}
  \exp\!\Big[\tfrac{1}{2g^{2}}\sum_{m,\rho\sigma}
    \bigl(e^{if_{\rho\sigma}(m)}+e^{-if_{\rho\sigma}(m)}\bigr)\Big].
\label{eq:3.15}
\end{equation}
这是算符 $e^{i\thl{\mu\nu}}$ 的传播子；它只定义在欧几里得时空格点的格点 $n$ 上。
如果格点间距为 $a$，这意味着传播子只对形如 $in_{0}a$ 的虚时间有定义，其中 $n_{0}$
是整数。

只在离散虚时间上有定义的理论远不能令人满意。幸运的是，可以推广理论来为量子化
理论定义哈密顿量。原则上，通过对角化哈密顿量可以得到理论的粒子本征态和散射振幅。
哈密顿量将用转移矩阵形式化来定义。这里只给出转移矩阵方法的简要讨论。关于这些
思想的综述${}^{6}$见 Wilson 和 Kogut。${}^{6}$ 在一个简单标量场理论中，包括单粒子
能量和散射振幅的近似计算的详细讨论由 Suri 给出。${}^{5}$

考虑 $Z$ 的表达式。对格点坐标 $n_{0}$ 引入有限界限，比如
\begin{equation}
-N \le n_{0} \le N.
\label{eq:3.16}
\end{equation}
引入周期性边界条件（见下文${}^{6}$）。于是可以把 $Z$ 写成矩阵 $V$ 的迹，更准确
地说
\begin{equation}
Z = \Tr V^{2N+1}.
\label{eq:3.17}
\end{equation}
这一公式之所以可能，是因为作用量 $A$ 中每一项涉及不超过两个相邻的 $n_{0}$ 值。

要建立矩阵 $V$，必须首先理解它所作用的函数空间。这里使用的空间是所有函数
$\Phi(\thl{ij})$ 的空间（对每个 $\thl{ij}$ 以 $2\pi$ 为周期），其中指标 $i$
只取 1 到 3，格点变量 $n$ 只有三个分量 $(n_{1},n_{2},n_{3})$。矩阵 $V$ 将定义为
两组自变量的函数，即 $\thl{ij}$ 和 $\thl{ij}'$，这两组自变量对应于相邻两个
$n_{0}$ 值的时空场 $\thl{ij}$。两个 $V$ 的矩阵乘法涉及对一组变量 $\thl{ij}$
的积分。把 $V$ 定义为
\begin{equation}
V = \prod_{n}\prod_{i,j}\int\ddth{i0}\,e^{U}.
\label{eq:3.18}
\end{equation}

下面的量 $U$ 写出来看起来很复杂，但它的全部内容就是作用量 $A$ 中涉及给定的
$n_{0}$ 最近邻对的那部分；$A$ 中只涉及单个 $n_{0}$ 值的项在连接 $n_{0}$ 与
$n_{0}+1$ 及 $n_{0}$ 与 $n_{0}-1$ 的矩阵之间平分。结果为
\begin{equation}
U = \frac{1}{2g^{2}}\sum_{n,i,j}\Bigl[
    \bigl(e^{if_{0i,j}(n)}+e^{-if_{0i,j}(n)}\bigr)
    +\bigl(e^{i\thl{i0,j}(n)}+e^{-i\thl{i0,j}(n)}\bigr)\Bigr],
\label{eq:3.19}
\end{equation}
其中
\begin{equation}
f_{0i}(n) = \thl{0i}(n)+\thl{i0}(n+\hat0)-\thl{0i}(n+\hat i)-\thl{i0}(n),
\label{eq:3.20}
\end{equation}
以及
\begin{equation}
f_{i0}(n) = \thl{i0}(n)+\thl{0i}(n+\hat i)-\thl{i0}(n+\hat0)-\thl{0i}(n).
\label{eq:3.21}
\end{equation}

根据这里给出的 $V$ 定义，容易看出迹 $\Tr V^{2N+1}$ 重现了 $Z$ 方程所涉及的
全部积分，而 $2N+1$ 个指数 $U$ 之和重现了作用量 $A$，只是额外增加了一些把边界
$n_{0}=N$ 与边界 $n_{0}=-N$ 耦合起来以实现周期结构的项。

注意 $\thl{\mu\nu}$ 的 0 分量受到了特殊处理。这是因为 $A$ 中没有涉及多个
$n_{0}$ 值的 $\thl{0\mu}$ 项；这使得可以把对 $\thl{0\mu}$ 的积分包含在 $V$ 的
定义中，而不是包含在矩阵乘法的定义中。

矩阵 $V$ 用于定义量子化理论。简要地说如下完成。$V$ 是厄米矩阵，即
\begin{equation}
V(\theta,\theta')^{*} = V(\theta',\theta).
\label{eq:3.22}
\end{equation}
[仔细检查式~\eqref{eq:3.18}--\eqref{eq:3.21} 可以验证这一结果。必须记住
$\thl{0\mu}$ 变量在 $V$ 的定义中被积掉了。这意味着在取 $V$ 的复共轭时，还可以
作变量替换 $\thl{0\mu} \to -\thl{0\mu}$。] 因此 $V$ 具有一整套正交完备的本征态
$\psi_{\lambda}$ 和本征值 $\lambda$。现在如下定义哈密顿量 $H$：$H$ 的本征态就是
$V$ 的本征态；$H$ 的本征值由下式给出
\begin{equation}
E_{\lambda} = -a^{-1}\ln\lambda,
\label{eq:3.23}
\end{equation}
其中 $\lambda$ 是 $V$ 的相应本征值。因子 $a^{-1}$ 的原因稍后就会清楚。使用对数的
原因是使多粒子散射态的能量成为单粒子能量之和（见 Suri${}^{5}$）。

如果 $V$ 有负本征值 $\lambda$，这个定义就会出现问题。如果出现这种情况，$H$ 将有
复本征值。Suri${}^{5}$ 研究的情形没有出现这种情况；这里是否会出现作者不知道；
这个问题必须进一步研究。即使 $V$ 有负本征值，它们在 $a\to0$ 极限下（如果该极限
存在）可能也是无关紧要的。

定义式~\eqref{eq:3.23} 意味着
\begin{equation}
V = e^{-aH}.
\label{eq:3.24}
\end{equation}
这意味着 $V$ 是把一个态经过虚时间 $ia$ 传播的算符。由此推得传播子 $D_{\mu\nu}$
是哈密顿量 $H$ 理论中虚时间 $in_{0}a$ 处的真空期望值。这一断言的证明见文献
5 和 6。

格点量子化程序可以推广到非阿贝尔规范理论。做法如下。用一组变量
$\thl{\mu\nu}^{a}$ 代替单一变量 $\thl{\mu\nu}$，其中 $a$ 是内指标。对每个 $n$
和 $\mu$，$\thl{\mu\nu}$ 用于参数化规范群的一个元素 $h_{\mu\nu}$。用
$U(\thl{\mu\nu})$ 代替 $\exp(i\thl{\mu\nu})$，其中 $U$ 是 $\thl{\mu\nu}$ 在夸克
表示中的酉矩阵。乘积 $\bar\psi_{n}\gamma_{\mu}\psi_{n+\hat\mu}\exp(i\thl{\mu\nu})$
替换为 $\bar\psi_{n}\gamma_{\mu}U(\thl{\mu\nu})\psi_{n+\hat\mu}$。规范变换由一组
群变换 $\gamma_{n}$ 定义；在这些变换下
\begin{equation}
\thl{\mu\nu} \to \thl{\mu\nu}^{a},
\label{eq:3.25}
\end{equation}
\begin{equation}
h_{\mu\nu} = \exp\!\Bigl(i\sum_{a}T^{a}\thl{\mu\nu}^{a}\Bigr),
\label{eq:3.26}
\end{equation}
\begin{equation}
\psi_{n} \to U(\gamma_{n})\psi_{n},\qquad
\bar\psi_{n} \to \bar\psi_{n}U(\gamma_{n})^{-1},
\label{eq:3.27}
\end{equation}
\begin{equation}
U(\thl{\mu\nu}) \to U(\gamma_{n})\,U(\thl{\mu\nu})\,U(\gamma_{n+\hat\mu})^{-1},
\label{eq:3.28}
\end{equation}
其中 $\gamma_{n}\,\thl{\mu\nu}\,\gamma_{n+\hat\mu}^{-1}$ 按群的乘法法则计算。
规范场的简单规范不变作用量是
\begin{equation}
\begin{split}
A_{G} = \frac{1}{2g^{2}}\sum_{n,\mu\nu}\Bigl[1-\tfrac{1}{2}\Tr
    &\bigl(U(\thl{\mu\nu}(n))\,U(\thl{\nu\mu}(n+\hat\mu))\\
    &\qquad U^{-1}(\thl{\mu\nu}(n+\hat\nu))\,U^{-1}(\thl{\nu\mu}(n))\bigr)\Bigr],
\end{split}
\label{eq:3.29}
\end{equation}
其中酉变换 $U_{A}$ 取自群的伴随表示。（任何表示都一样适用。）在经典连续极限下，
$\thl{\mu\nu}$ 变成无穷小群变换 $\thl{\mu\nu} = ga\,A_{\mu}^{a}(na)T^{a}$，其中
$A_{\mu}^{a}$ 在 $a\to0$ 时固定；在该极限下可以证明（需费些力气）$A_{G}$ 化为
标准的连续杨-米尔斯作用量。对非阿贝尔理论量子化时，对 $\thl{\mu\nu}$ 的积分是
对群的所有元素作的群积分。这些是紧致积分（对紧致群而言），因此无需规范固定项。

由于用阿贝尔理论即可说明夸克束缚，这里不再进一步研究非阿贝尔理论。

最后讨论狄拉克场的量子化。方便的做法是先用与规范场类似的方式量子化狄拉克场。
唯一出现的问题是定义费米场的``积分''。这是可以做到的。${}^{11}$

对量子化至关重要的积分性质是积分变量的平移不变性。例如，在格点上量子化标量场
$\phi$ 时，场平均涉及积分 $\int\dd\phi_{n}$，后者具有平移不变性
\begin{equation}
\int\dd\phi_{n}\,f(\phi_{n}+\xi_{n}) = \int\dd\phi_{n}\,f(\phi_{n})
\label{eq:3.30}
\end{equation}
对任意可积函数 $f$ 和任意常数 $\xi_{n}$ 成立。正是这种平移不变性使费曼路径积分
成为施温格作用量原理的实现（进一步讨论见文献 12）。类似地，需要为费米场定义具有
相同平移不变性的积分。抽象地说，要定义一个括号运算 $(\ )$，它作用在纯粹反交换
的费米场 $\psi_{n}$ 的函数上，具有性质
\begin{equation}
(f) = \int f(\psi,\bar\psi)\,\prod_{n}\dd\psi_{n}\,\dd\bar\psi_{n},
\label{eq:3.31}
\end{equation}
其中 $\psi_{n}$ 和 $\bar\psi_{n}$ 是反交换的 $c$ 数（它们由施温格引入）。括号运算
应当对每个函数 $f$ 产生一个数；它还应当是线性运算。因此对有限格点，只需为
$\psi_{n}$ 和 $\bar\psi_{n}$ 的所有单项式规定括号 $(\ )$。由于反交换规则，
$\psi_{n}^{2}$ 和 $\bar\psi_{n}^{2}$ 为零（更准确地说，$\psi_{n\alpha}^{2}$ 和
$\bar\psi_{n\alpha}^{2}$ 为零，其中 $\psi_{n\alpha}$ 和 $\bar\psi_{n\alpha}$ 是
$\psi_{n}$ 与 $\bar\psi_{n}$ 的任一分量），因此可能的单项式只有有限个。现在容易
看出，括号 $(\ )$ 对所有 $\psi$ 与 $\bar\psi$ 的乘积都必须为零，除非该乘积包含
所有可能的 $\psi$ 和 $\bar\psi$。例如，设有两个格点 $0$ 和 $1$，且 $\psi_{n}$ 与
$\bar\psi_{n}$ 是单分量场。那么括号必须满足
\begin{equation}
(\psi_{0}\bar\psi_{0}\psi_{1}\bar\psi_{1}) = 1,
\label{eq:3.32}
\end{equation}
其中 1 是任意选取的常数。括号运算的这一定义满足平移不变性：例如，因为乘上
$\psi$ 的项全为零。还要注意反交换规则意味着，例如
\begin{equation}
(\bar\psi_{0}\psi_{0}\bar\psi_{1}\psi_{1}) = -1.
\label{eq:3.33}
\end{equation}
（与标量情形类比，要求 $\psi\bar\psi = 1$，而不是 $\psi\bar\psi = 1 - \bar\psi\psi$。）

现在可以定义包含狄拉克场的完整规范理论在格点上的费曼路径积分。例如，格点上的
流-流传播子为
\begin{equation}
D_{\mu\nu}(x) = Z^{-1}\prod_{n,\rho}\int\dd\psi_{n}\,\dd\bar\psi_{n}\,
  \dd\thl{\rho}\,J_{\mu}(x)J_{\nu}(0)\,e^{A},
\label{eq:3.34}
\end{equation}
其中 $A$ 是式~\eqref{eq:3.12} 的完整作用量，而
\begin{equation}
J_{\mu}(n) = \bar\psi_{n}\gamma_{\mu}\psi_{n}.
\label{eq:3.35}
\end{equation}

这一路径积分形式不同于第二节讨论的形式。不过，从现在的表达式很容易导出第二节
路径积分的格点形式。做法是把式~\eqref{eq:3.34} 按 $\kappa$ 展开，其中 $\kappa$
是最近邻耦合项 $\bar\psi_{n}\gamma_{\mu}\psi_{n+\hat\mu}e^{i\thl{\mu\nu}}$ 等的
系数。这一最近邻耦合项可以用从格点 $n$ 到格点 $n+\hat\mu$ 的一条线在图上表示。
最好通过研究式~\eqref{eq:3.34} 分子展开中的一项的例子来描述这一展开，现在就来
讨论它。展开项的一个例子示意性地画在图~\ref{fig:9} 中。该项的表达式为
\begin{equation}
\begin{split}
D = {} & (-1)\kappa^{4}\,
    e^{i[\thl{\mu\nu}(0,0)+\thl{\mu\nu}(1,0)+\thl{\mu\nu}(0,1)+\thl{\mu\nu}(1,1)]}\\
      & \times \bigl(\bar\psi_{00}\gamma_{\mu}\psi_{10}\bigr)
        \bigl(\bar\psi_{10}\gamma_{\nu}\psi_{11}\bigr)
        \bigl(\bar\psi_{11}\gamma_{\mu}\psi_{01}\bigr)
        \bigl(\bar\psi_{01}\gamma_{\nu}\psi_{00}\bigr),
\end{split}
\label{eq:3.36}
\end{equation}
（$D$ 表示 diagram），其中涉及的四个格点是 $(n_{0},n_{1}) = (0,0)$、$(1,0)$、
$(0,1)$ 和 $(1,1)$；$n_{2}$ 和 $n_{3}$ 的值不变，在式~\eqref{eq:3.36} 的记法中
被省略。作用量 $A_{0}$ 略去 $\kappa$ 项，即
\begin{equation}
A_{0} = -\epsilon\sum_{n}\bar\psi_{n}\psi_{n}
  + \frac{1}{2g^{2}}\sum_{n,\mu\nu}
    \bigl(e^{if_{\mu\nu}(n)}+e^{-if_{\mu\nu}(n)}\bigr).
\label{eq:3.37}
\end{equation}

\begin{figure}[t]
  \centering
  \includegraphics[width=0.3\textwidth]{fig9}
  \refstepcounter{figure}%
  \caption*{图 9. 格点上的基本小方块。}
  \label{fig:9}
\end{figure}

$D$ 的计算分两部分，一是对所有规范场 $\thl{\mu\nu}$ 积分，二是计算 $\psi,\bar\psi$
括号。这两个计算相互独立，即 $D$ 分解为 $D_{\theta}D_{\psi}$。量 $D_{\theta}$ 为
\begin{equation}
D_{\theta} = \prod_{n,\mu\nu}\int\ddth{\mu\nu}
  \exp\!\Bigl[i\bigl(\thl{\mu\nu}(0,0)+\thl{\mu\nu}(1,0)
    -\thl{\mu\nu}(1,1)-\thl{\mu\nu}(0,1)\bigr)+A_{0}\Bigr].
\label{eq:3.38}
\end{equation}

这是闭合圈上线积分的指数之规范场平均的一个例子，该圈即图~\ref{fig:9} 中的圈。
$\psi$ 括号计算可以进一步分解为对每个格点独立的括号计算，因为 $A_{0}$ 中不包含
耦合不同格点的 $\psi$ 与 $\bar\psi$ 项。为简单起见，只考虑圈上的四个格点。利用
反交换规则稍微移动 $\psi$ 之后，括号变为
\[
\begin{split}
&(\bar\psi_{00}e^{-i\thl{\mu\nu}(0,0)}\psi_{00})
 (\bar\psi_{10}e^{-i\thl{\mu\nu}(1,0)}\psi_{10})\\
&\qquad\times
 (\bar\psi_{11}e^{-i\thl{\mu\nu}(1,1)}\psi_{11})
 (\bar\psi_{01}e^{-i\thl{\mu\nu}(0,1)}\psi_{01}).
\end{split}
\]
要得到所有可能的 $\psi$ 与 $\bar\psi$ 的乘积，就意味着必须取所有可能的
$\psi_{00}$ 与 $\bar\psi_{00}$、所有可能的 $\psi_{10}$ 与 $\bar\psi_{10}$ 等等的
乘积。总之，完整括号可以写成四个独立括号的乘积。定义
\begin{equation}
D_{\psi,0} = \bigl(\bar\psi_{00}e^{i\thl{\mu\nu}(0,0)}\psi_{10}\bigr),
\label{eq:3.39}
\end{equation}
\begin{equation}
D_{\theta,0} = \bigl(\bar\psi_{10}\psi_{10}\bigr)_{0}.
\label{eq:3.40}
\end{equation}
由于 $\psi_{10}$、$\bar\psi_{10}$、$\psi_{00}$、$\bar\psi_{00}$ 以及 $\gamma$
矩阵所隐含的自旋指标，$D_{\theta}$ 和 $D_{\psi}$ 都是自旋空间中的矩阵。完整括号
就是
\begin{equation}
D_{\psi} = -\kappa^{4}\Tr\bigl(D_{\psi,0}\gamma_{\mu}
  D_{\psi,0}\gamma_{\nu}D_{\psi,0}\gamma_{\mu}D_{\psi,0}\gamma_{\nu}\bigr).
\label{eq:3.41}
\end{equation}

矩阵 $D_{\psi,0}$ 和 $D_{\theta,0}$ 容易确定。例如，$D_{\theta,0}$ 的显式形式为
\begin{equation}
D_{\theta,0} = \delta_{\alpha\beta}
  - \epsilon\bigl(\gamma_{0}\gamma_{0}\bigr)_{\alpha\beta}
  - \epsilon\bigl(\gamma_{1}\gamma_{1}\bigr)_{\alpha\beta}
  - \epsilon\bigl(\gamma_{2}\gamma_{2}\bigr)_{\alpha\beta}
  - \epsilon\bigl(\gamma_{3}\gamma_{3}\bigr)_{\alpha\beta}.
\label{eq:3.42}
\end{equation}
指数可以按 $\epsilon$ 的幂展开；假定旋量有四个分量，只有 $\epsilon^{4}$ 项才能
产生全部四个 $\psi_{10}$ 乘全部四个 $\bar\psi_{10}$ 的乘积；结果是
\begin{equation}
D_{\theta,0} = \epsilon^{4}\bigl(\gamma_{0}\gamma_{1}\gamma_{2}\gamma_{3}\bigr)_{\alpha\beta},
\label{eq:3.43}
\end{equation}
（负号来自如下约定：当 $\bar\psi_{10}$ 出现在 $\psi_{10}$ 左边时括号为正）。
类似的计算给出
\begin{equation}
D_{\psi,0} = \bigl(\psi_{10}\bar\psi_{10}\bigr)_{0},
\label{eq:3.44}
\end{equation}

这一例子的结果很容易推广。一般阶数 $\kappa^{m}$ 的项只有当最近邻耦合组合成闭合圈
时才非零（开线端点的格点会多出 $\psi_{n}$ 或 $\bar\psi_{n}$，因此该格点的括号
为 0）。闭合圈的括号计算给出一个迹，其中圈中每条线对应 $\kappa$ 乘以一个
$\gamma$ 矩阵，圈中每个格点对应一个 $D$（除去有流出现的点 $n$ 和 $0$）。
对规范场的平均涉及每个圈周围 $\thl{\mu\nu}$ 之和的指数。圈数可以任意多。
